% Revised methodology draft for the SHIFT-ATCP manuscript.
% All numerical results must be regenerated from the revised pipeline.

\section{Study Design and Cohort Harmonization}
\label{sec:study_design}

We reformulate the study as a cross-survey transportability analysis rather than
using a source cohort with undocumented clinical provenance. The primary source
cohort is the 2023 Behavioral Risk Factor Surveillance System (BRFSS), and the
external target cohort is the 2017--2018 National Health and Nutrition
Examination Survey (NHANES). BRFSS is a large, repeated cross-sectional
telephone survey; the CDC reports 433,323 records for the 2023 annual file \cite{CDC_BRFSS_2023}.
The raw survey files, codebooks, release dates, and preprocessing scripts are
retained to make the cohort construction auditable. The Pima Indians Diabetes
Dataset is retained only as a secondary, small-sample stress test and is not
used to support population-representative claims.

The primary outcome is \emph{self-reported clinician-diagnosed diabetes}.
BRFSS respondents coded as having been told they have diabetes are positive;
respondents reporting no diabetes are negative. Pregnancy-only diabetes,
pre-diabetes/borderline responses, refusal, and unknown responses are excluded.
NHANES uses the corresponding self-reported diagnosis item (DIQ010), with
positive and negative responses retained and borderline, unknown, refused, and
missing responses excluded. This outcome does not represent all diabetes,
because undiagnosed laboratory-defined diabetes is not captured. A separate
laboratory-defined outcome analysis, if conducted, must exclude the diagnostic
laboratory measurements from the predictors.

Both cohorts are restricted to adults aged 20 years or older to align the
NHANES adult questionnaire variables. The prespecified common predictor set is
age, sex (female indicator), body-mass index, self-reported hypertension,
heart-disease history, and ever-smoking history. The BRFSS BMI variable is
converted from its stored two-decimal representation to kg/m$^2$. Hypertension
is positive only for an affirmative non-pregnancy response; negative responses
are coded zero, while pregnancy-only, borderline/pre-hypertension, unknown,
and refused responses are treated as missing. Heart-disease history is positive
if either prior myocardial infarction or coronary heart disease/angina is
reported, negative only if both component questions are negative, and otherwise
missing. Smoking history indicates having smoked at least 100 cigarettes.
Definitions and missing-value rules are fixed before model fitting and applied
identically in both cohorts \cite{CDC_BRFSS_CV_2023,NCHS_NHANES_2017_2018}. No predictor is created as an all-missing column
in either cohort.

BRFSS design variables (_LLCPWT, _STSTR, _PSU) and NHANES interview weights,
strata, and primary sampling units (WTINT2YR, SDMVSTRA, SDMVPSU) are retained.
The primary machine-learning and conformal analysis is an individual-level
analytic-sample analysis. Survey-weighted test metrics and PSU-within-stratum
cluster-bootstrap intervals are reported as a sensitivity analysis. Because
the primary random split is at the respondent level, this cluster bootstrap is
an approximation rather than a complete design-based variance estimator.
Moreover, survey weights are not incorporated into model fitting,
density-ratio estimation, or conformal calibration. Weighted summaries are
therefore not interpreted as a full design-based conformal coverage guarantee.

\section{Leakage-Controlled Experimental Protocol}
\label{sec:protocol}

For each random seed, the source cohort is split into 60\% model training,
10\% probability-calibration fitting, and 30\% source conformal calibration.
The source probability calibrator is fitted only on the second partition, and
source conformal scores are computed only on the third partition. The target
cohort is randomly partitioned without using its labels to allocate rows:
15\% unlabeled covariate adaptation, 15\% labelled method-gating, 20\%
target-model training, 10\% target-model probability calibration, 20\%
target conformal calibration, and 20\% final test. The target adaptation
partition is used only for covariate-shift estimation. Target-test covariates are
queried only to obtain their pointwise density ratios (a transductive use of
unlabelled test features); target-test labels remain inaccessible until final
evaluation. Target labels are used
only in the explicitly identified gate, target-trained reference, calibration,
and conformal-calibration partitions. The final test labels are not used for
fitting, calibration, gating, hyperparameter selection, or method selection.

A source-trained classifier $f_\theta$ is fitted on source training data. Its
probability output is calibrated by an independent Platt map $g_\phi$ fitted on
the source probability-calibration partition:
\begin{equation}
  \widehat p(x)=g_\phi(f_\theta(x)).
\end{equation}
For a labelled calibration observation $(x_i,y_i)$, the nonconformity score is
\begin{equation}
  s_i=1-\widehat p_{y_i}(x_i),\qquad
  \widehat p_1(x)=\widehat p(x),\quad
  \widehat p_0(x)=1-\widehat p(x).
\end{equation}
For $n$ calibration scores, the ordinary split-conformal threshold is the
$\lceil(n+1)(1-\alpha)\rceil$-th order statistic, with threshold $+\infty$ if
the requested order statistic exceeds the available calibration sample.

\section{Shift Diagnostics and Importance-Weighted Baseline}
\label{sec:shift_diagnostics}

Shift is quantified on the same prespecified common predictors in both cohorts.
For each predictor, we report source and target missingness separately and
calculate the population stability index (PSI) on shared bins:
\begin{equation}
  \mathrm{PSI}=\sum_b (p_{T,b}-p_{S,b})\log\frac{p_{T,b}}{p_{S,b}}
  =D_{\mathrm{KL}}(p_T\Vert p_S)+D_{\mathrm{KL}}(p_S\Vert p_T).
\end{equation}
We verify this identity numerically and report both KL terms, rather than
presenting incompatible values. A cross-validated domain classifier is trained
on observed values of the common predictors; missingness-rate differences are
reported separately so that missing-feature artifacts are not silently
interpreted as biological covariate shift. The out-of-fold domain AUC and
importance-weight effective sample size are reported for every run.

As a comparator, we estimate the density ratio
$r(x)=p_T(x)/p_S(x)$ using a cross-fitted domain classifier. If
$e(x)=P(D=T\mid X=x)$ is estimated by a domain classifier with balanced
domain-class weights, then
\begin{equation}
 \widehat r(x)=\frac{\widehat e(x)}{1-\widehat e(x)}.
\end{equation}
The domain classifier is fitted with balanced domain-class weights, so that its
effective source and target priors are equal; the prior-odds correction is
therefore one.
Weights are clipped only at prespecified bounds, and clipping fractions and
weight effective sample size are reported. For a test point $x$, the
importance-weighted source threshold is
\begin{equation}
 \widehat q_{IW}(x)=\inf\left\{q:
 \sum_{i\in S_{cal}}\widehat r(x_i)\,\mathbb{I}(s_i\le q)
 \ge (1-\alpha)\left[\sum_{i\in S_{cal}}\widehat r(x_i)+\widehat r(x)\right]
 \right\}.
\end{equation}
If the required mass exceeds the observed calibration mass, the threshold is
$+\infty$. The test-point weight is therefore included explicitly. This
construction has a finite-sample weighted-conformal interpretation when the
true density ratios are known and the covariate-shift assumption
$P_T(Y\mid X)=P_S(Y\mid X)$ holds. In this study the ratios are estimated, so
we describe the resulting method as an approximate importance-weighted
baseline and do not claim an exact finite-sample coverage guarantee.

\section{Calibration-Gated SHIFT-ATCP}
\label{sec:shift_atcp}

The revised SHIFT-ATCP is a calibration-gated transport procedure, not the
previous scalar source-trust formula. It compares the importance-weighted
source threshold with target-only split conformal using a labelled target gate
partition that is disjoint from target conformal calibration and final test.
Let $G$ denote the gate set. The weighted-source candidate is accepted only if
(i) its one-sided 90\% Wilson lower confidence bound for gate coverage is at
least $1-\alpha-\delta$, and (ii) its gate mean prediction-set size is at
least $\eta$ smaller than that of the target-only marginal conformal baseline.
The constants $\delta$ and $\eta$ are prespecified and varied in sensitivity
analysis. Otherwise, SHIFT-ATCP falls back to the target-only threshold:
\begin{equation}
 \widehat q_{SHIFT}(x)=
 \begin{cases}
   \widehat q_{IW}(x), & \mathrm{Gate}(G)=1,\\
   \widehat q_T, & \mathrm{Gate}(G)=0.
 \end{cases}
\end{equation}
The gate is evaluated only on its designated target partition. It is not
re-evaluated or modified using final test labels. A high fallback frequency is
reported as evidence that source transport was not supported by the available
calibration data, rather than being presented as a successful source-transfer
result. The method is not expected to outperform target-only calibration in
every cohort; any claimed benefit must be supported by the complete external
results grid and controlled shift experiments.

\section{Class-Conditional Conformal Baseline and Target-Trained Reference}
\label{sec:mondrian}

To prevent marginal coverage from concealing poor diabetes-positive coverage,
we include ordinary target-only Mondrian conformal prediction as a primary
baseline. Separate target calibration quantiles are computed within each true
class:
\begin{equation}
 q_{T,k}=\mathrm{Quantile}_{1-\alpha}
 \{1-\widehat p_k(x_i): i\in T_{conf}, y_i=k\},\quad k\in\{0,1\}.
\end{equation}
The prediction set is
\begin{equation}
 \Gamma(x)=\{k\in\{0,1\}:1-\widehat p_k(x)\le q_{T,k}\}.
\end{equation}
This is the direct class-conditional baseline and replaces heuristic claims
based solely on a custom class-conditional selection score. We additionally
train a target-supervised reference classifier on the separate target training
partition, calibrate it on target calibration-fit data, and conformal-calibrate
it on target conformal-calibration data. This reference quantifies the
performance gap between source transport and a model trained with target labels;
it is not presented as a label-free transport method.

\section{Evaluation, Uncertainty, and Sensitivity Analysis}
\label{sec:evaluation}

Every model, method, seed, and cost vector is retained in the results grid.
No best model, feature set, or scenario is selected by final target-test
metrics. We report marginal coverage, class-0 coverage, diabetes-positive
coverage, absolute class-coverage gap, mean set size, singleton rate, empty-set
rate, referral rate, false-positive and false-negative singleton rates, Brier
score, AUROC, and 10-bin expected calibration error. Clinical referral cost is
computed from an explicit action rule: singleton sets trigger the corresponding
class prediction, while empty or two-class sets are referred. We vary the
false-positive, false-negative, and referral costs over prespecified vectors.

For each seed and method, percentile bootstrap intervals are computed for
marginal and class-specific coverage, class-coverage gap, set size, referral
rate, and clinical cost. When survey design variables are supplied, the
sensitivity analysis resamples primary sampling units within strata; otherwise
participant-row bootstrap intervals are reported. Repeated random seeds are
used to quantify split variability, and all seed-level results are released.
The main tables retain model and configuration identifiers rather than
averaging across unrelated models, feature sets, or scenarios.

Additional analyses include calibration-sample-size sensitivity, density-ratio
clipping and effective-sample-size sensitivity, gate-tolerance sensitivity,
controlled covariate-shift experiments where $P(Y\mid X)$ is fixed, and
prevalence/label-shift experiments where the covariate-shift assumption is
violated. Pima results are reported separately as a small, female-only
stress-test cohort. Population-level interpretations are avoided unless the
survey design is handled consistently throughout the corresponding analysis.
